\section{Conclusion and Future Work}
\label{sec:conclusion}

In this paper, we have explored the vulnerabilities of current test-time dynamic model merging pipelines through a strictly empirical and statistical lens. We exposed two major weaknesses of existing architectures: the sensitivity of quantum-inspired wave-superposition models (QWS-Merge) under low-data, sample-wise vectorized calibration splits, and the presence of the heterogeneity collapse confounder where mixed-task inference streams dilute sample-specific routing benefits.

Our deconstruction has led to a profound demystification of dynamic model merging. Through a rigorous baseline audit, we compared standard dynamic routers against our well-regularized Softmax baseline (\texttt{L3\_Softmax\_WellReg}). Our findings conclusively prove that proper prior-layer design---specifically, zero-initialized Softmax routing layers regularized with $L_2$ weight decay---completely resolves Vectorization Collapse and mitigates heterogeneity collapse without requiring any custom or explicit training losses. An exhaustive empirical audit across 10 independent random seeds, regularizing parameter sweeps, and batch heterogeneity stress tests demonstrated that our well-regularized baseline maintains high, flatline stability across all deployment batch sizes ($B=1$ to $B=512$), outperforming all standard unregularized dynamic baselines.

Crucially, our quantitative analysis of learned weights reveals that the routing coefficients stay remarkably close to their maximum-entropy uniform prior, yielding a Mean Absolute Deviation of only \textbf{0.0236} (or \textbf{2.36\%}) from $0.25$. This low deviation exposes the \emph{Dynamic Routing Paradox}: to remain stable under data scarcity, the router must be so heavily regularized that its functional flexibility is severely limited. Consequently, the well-regularized router only improves joint accuracy by \textbf{1.16\%} over static Uniform Merging (59.16\% vs. 58.00\%). Given that Uniform Merging requires zero calibration data, zero optimization time, and zero runtime memory or compute overhead (compared to the $O(B \cdot M)$ parameter footprint expansion and reduced GPU arithmetic intensity of real-time parameter assembly), we advocate for naive Uniform Merging as an exceptionally strong, cost-free baseline that dynamic designs must beat.

Our findings advocate for a fundamental shift in the design of dynamic routers for model merging: rather than adopting convoluted, non-monotonic mathematical metaphors that complicate the optimization landscape or relying on redundant, explicit training losses, researchers should prioritize simple classical linear projections regularized with appropriate architectural priors. Proper zero-initialization and weight decay are the most critical, yet under-appreciated, design elements for stable test-time model merging.
